% =====================================================================
%  CAPÍTULO 3. PRESENTACIÓN DE RESULTADOS Y ANÁLISIS DE LOS MISMOS
% =====================================================================
\chapter{PRESENTACIÓN DE RESULTADOS Y ANÁLISIS DE LOS MISMOS}
\label{cap:resultados}

\guia{Según el Taller 5, los resultados se presentan en tablas y gráficos fáciles de entender. Se describen en pasado, en párrafos y sin juicios de valor, y responden a cada objetivo con la metodología de análisis definida. La interpretación y el contraste con la teoría y los antecedentes van en la sección de discusión.}

% ---------------------------------------------------------------------
\section{Resultados del diagnóstico y de la línea base}
\guia{OE1 y OE2. Porcentaje de cumplimiento de la norma ISO~9001:2026 por cláusula (gráfico de barras); no conformidades por categoría (Pareto); diagrama de Ishikawa; tabla de índices de desperdicio y costo de no calidad de la primera losa por material.}

% ---------------------------------------------------------------------
\section{Resultados de la aplicación del SGC en el control de materiales}
\guia{OE4. Índice de desperdicio por material y por paño en la segunda losa, comparado con la línea base (tabla y gráfico de barras agrupadas). Porcentaje de lotes conformes y de paños liberados sin observaciones.}

% ---------------------------------------------------------------------
\section{Resultados del control de calidad de materiales y del hormigón}
\guia{Resultados de asentamiento y de resistencia a compresión a 28 días de la losa: media, desviación estándar, coeficiente de variación y cumplimiento de los criterios de aceptación (sección~\ref{sec:contraste} y Tabla~\ref{tab:ppi}). Comparar con los resultados de la línea base (Tabla~\ref{tab:probetas}).}

% ---------------------------------------------------------------------
\section{Resultados de costos}
\guia{Costo unitario con y sin SGC a precios constantes (ecuación~\ref{eq:cu}), comparado con el precio contractual de Bs~215,68/\mdos{}. Costos de la calidad y de no calidad (ecuación~\ref{eq:cnc}). Ahorro proyectado en el saldo de 2.234,60~\mdos{}.}

% ---------------------------------------------------------------------
\section{Contraste de hipótesis}
\guia{Para cada hipótesis presentar: estadístico de prueba, valor $p$, decisión (se rechaza o no $H_0$) y conclusión en palabras. Incluir la prueba de normalidad.}

% ---------------------------------------------------------------------
\section{Validación del SGC por juicio de expertos}
\guia{Tabla con el coeficiente V de Aiken por criterio (pertinencia, claridad, suficiencia y aplicabilidad) y las observaciones de los expertos incorporadas al SGC.}

% ---------------------------------------------------------------------
\section{Discusión de resultados}
\guia{Contrastar los resultados con los antecedentes y el marco teórico. ¿Los índices de desperdicio y la reducción de costos son coherentes con lo reportado en otras investigaciones? Explicar las coincidencias y las diferencias, las limitaciones del estudio y el alcance de la validez de los resultados.}
